% from ega3-1-fr.tex
\oldpage[0\textsubscript{III}]{53}
qui proviennent de l'homomorphisme $\psi^*(\mathcal F)\to f^*(\mathcal F)$
et du caractère fonctoriel des homomorphismes canoniques de (G, II, 5.2),
pour obtenir la commutativité de
(\hyperref[0.12.1.4.2-fr]{12.1.4.2}).

On notera que si $A=\Gamma(X,\mathcal O_X)$,
$A'=\Gamma(X',\mathcal O_{X'})$, l'homomorphisme (12.1.4.3) est un
di-homomorphisme de modules
correspondant aux anneaux $A$ et $A'$. On a une propriété de
\emph{transitivité} de (\hyperref[0.12.1.4.1-fr]{12.1.4.1}) pour le composé de
deux morphismes, analogue à la transitivité de
(\hyperref[0.12.1.3.5-fr]{12.1.3.5}). Enfin, notons que dans les définitions
précédentes, au lieu d'une résolution injective $\mathcal L^\bullet$ de
$\mathcal F$, on aurait aussi bien pu partir d'une résolution telle que
$H^p(X,\mathcal L^i)=0$ pour tout $i$ et tout $p>0$ (G, II, 4.7.1).
\end{env}

\begin{env}[12.1.5]
\label{0.12.1.5-fr}
Soient $\mathcal F$, $\mathcal G$, $\mathcal H$ trois
$\mathcal O_X$-Modules, et considérons un $\mathcal O_X$-homomorphisme
$u:\mathcal F\otimes_{\mathcal O_X}\mathcal G\to\mathcal H$, qui donne pour
la cohomologie des homomorphismes
\[
\label{0.12.1.5.1-fr}
\tag{12.1.5.1}
H^p(X,\mathcal F)\otimes_A H^q(X,\mathcal G)
\longrightarrow H^{p+q}(X,\mathcal H)
\]
déduits du cup-produit (\hyperref[0.12.1.2.1-fr]{12.1.2.1}). Montrons qu'avec
les hypothèses et notations de (\hyperref[0.12.1.3-fr]{12.1.3}), on a des
diagrammes commutatifs
\[
\label{0.12.1.5.2-fr}
\tag{12.1.5.2}
\xymatrix@C=10pt{
H^p(X,\mathcal F)\otimes_A H^q(X,\mathcal G)\ar[r]\ar[d]
&H^{p+q}(X,\mathcal H)\ar[d]\\
H^p(X',f^*(\mathcal F))\otimes_{A'}H^q(X',f^*(\mathcal G))\ar[r]
&H^{p+q}(X',f^*(\mathcal H))
}
\]
où les flèches verticales proviennent des homomorphismes canoniques
(\hyperref[0.12.1.3.5-fr]{12.1.3.5}). Pour cela, rappelons que
(\hyperref[0.12.1.5.1-fr]{12.1.5.1}) peut s'obtenir en partant des résolutions
\emph{canoniques} (G, II, 4.3) $\mathcal L^\bullet$, $\mathcal M^\bullet$,
$\mathcal N^\bullet$ de $\mathcal F$, $\mathcal G$, $\mathcal H$
respectivement (qui sont formées de $\mathcal O_X$-Modules), de l'application
linéaire $\mathcal L^\bullet\otimes_{\mathcal O_X}\mathcal M^\bullet
\to\mathcal N^\bullet$ de complexes de $\mathcal O_X$-Modules correspondant
à $u$, qui fournit un homomorphisme de complexes de $A$-modules
$\Gamma(X,\mathcal L^\bullet)\otimes_A\Gamma(X,\mathcal M^\bullet)
\to\Gamma(X,\mathcal N^\bullet)$, et en passant à la cohomologie des
homomorphismes
\[
H^p(\Gamma(X,\mathcal L^\bullet))\otimes_A
H^q(\Gamma(X,\mathcal M^\bullet))
\longrightarrow H^{p+q}(\Gamma(X,\mathcal N^\bullet))
\]
(G, II, 6.6). Or, on a évidemment un diagramme commutatif
\[
\label{0.12.1.5.3-fr}
\tag{12.1.5.3}
\xymatrix@C=8pt{
\Gamma(X,\mathcal L^\bullet)\otimes_A\Gamma(X,\mathcal M^\bullet)
\ar[r]\ar[d]
&\Gamma(X,\mathcal N^\bullet)\ar[d]\\
\Gamma(X',\psi^*(\mathcal L^\bullet))
\otimes_{\Gamma(X',\psi^*(\mathcal O_X))}
\Gamma(X',\psi^*(\mathcal M^\bullet))\ar[r]
&\Gamma(X',\psi^*(\mathcal N^\bullet))
}
\]
